\begin{resumo}[Abstract]
 \begin{otherlanguage*}{english}
This study presents the development of a chatbot aimed at providing academic support to students at the University of Brasília (UnB), based on a user-centered approach grounded in ITIL v4 and Design Thinking principles. The proposal stems from the need to centralize and simplify access to institutional information, which is often scattered and difficult to navigate. The solution consists of a functional bot integrated with widely used platforms such as WhatsApp, Telegram, and a complementary web application. The tool was developed by structuring official UnB documents into a database, enabling the chatbot to respond to frequently asked questions regarding deadlines, administrative procedures, courses, final papers (TCCs), internships, among other topics. To validate the solution’s effectiveness, a survey was conducted with students from the Gama campus (FCTE), whose responses directly informed the system’s requirements. Additionally, the chatbot provides a channel for redirecting users to human support when necessary, connecting them to institutional sectors such as the academic office. Based on the gathered requirements, a Minimum Viable Product (MVP) was implemented to demonstrate how the solution can enhance students' autonomy and efficiency in accessing information while also improving communication between students and the university.

   \vspace{\onelineskip}
 
   \noindent 
   \textbf{Key-words}: chatbot. academic support. ITIL v4. Design Thinking. University of Brasília.
 \end{otherlanguage*}
\end{resumo}
